\documentclass[UTF8,12pt]{ctexart}
\usepackage[a4paper,margin=24mm]{geometry}
\usepackage{amsmath,amssymb,bm,graphicx,adjustbox}
\newcommand{\fitmath}[1]{\begin{adjustbox}{max width=\linewidth}$\displaystyle #1$\end{adjustbox}}
\setlength{\parindent}{0pt}
\setlength{\parskip}{.7em}
\sloppy
\allowdisplaybreaks
\begin{document}
\section*{1994 年考研数学三 · 参考答案}
\subsection*{1994 年真题参考答案（试卷 IV）}

\subsection*{一、填空题}

（1）\(\displaystyle \ln3\)。 （2）1。 （3）\(\displaystyle -\dfrac{\displaystyle ye^{xy}+\sin x}{\displaystyle xe^{xy}+2y}\)。


（4）\(\displaystyle \begin{pmatrix}0&0&\cdots&0&\dfrac{\displaystyle 1}{\displaystyle a_n}\\\dfrac{\displaystyle 1}{\displaystyle a_1}&0&\cdots&0&0\\0&\dfrac{\displaystyle 1}{\displaystyle a_2}&\cdots&0&0\\\vdots&\vdots&&\vdots&\vdots\\0&0&\cdots&\dfrac{\displaystyle 1}{\displaystyle a_{n-1}}&0\end{pmatrix}\)。 （5）\(\displaystyle \dfrac{\displaystyle 9}{\displaystyle 64}\)。


\subsection*{二、选择题}

（1）B。 （2）C。 （3）C。 （4）D。 （5）B。


三、\(\displaystyle \dfrac{\displaystyle 3\pi}{\displaystyle 2}\)。


四、1。


五、\(\displaystyle \dfrac{\displaystyle \partial^2f}{\displaystyle \partial x\partial y}=\dfrac{\displaystyle x^2-y^2}{\displaystyle x^2+y^2}\)。


六、\(\displaystyle \lim_{x\to0}\dfrac{\displaystyle F(x)}{\displaystyle x^{2n}}=\dfrac{\displaystyle 1}{\displaystyle 2n}f^{\prime}(0)\)。


七、（1）\(\displaystyle a=\dfrac{\displaystyle 1}{\displaystyle e}\)，切点为 \(\displaystyle (e^2,\allowbreak 1)\)。（2）\(\displaystyle \dfrac{\displaystyle \pi}{\displaystyle 2}\)。


八、证明略。（证明 \(\displaystyle F^{\prime}(x)>0\)（\(\displaystyle x>a\)）。）


九、（1）证明略。（原线性方程组的增广矩阵的行列式是一个范德蒙德行列式。）（2）\(\displaystyle x=k(2,\allowbreak 0,\allowbreak -2)^{\mathrm T}+(-1,\allowbreak 1,\allowbreak 1)^{\mathrm T}\)，其中 \(\displaystyle k\) 为任意常数。


十、\(\displaystyle x+y=0\)。


十一、\(\displaystyle X\sim\begin{array}{c|ccc}X&-1&0&1\\P&0.1344&0.7312&0.1344\end{array}\)。


十二、当 \(\displaystyle \mu\approx10.9\) 毫米时，平均利润最大。


\subsection*{1994 年真题参考答案（试卷 V）}

\subsection*{一、填空题}

（1）【同试卷 IV 第一、（1）题】 （2）【同试卷 IV 第一、（2）题】 （3）【同试卷 IV 第一、（3）题】 （4）【同试卷 IV 第一、（4）题】 （5）\(\displaystyle \dfrac{\displaystyle 2}{\displaystyle 3}\)。


\subsection*{试卷 V（续）}

\subsection*{二、选择题}

（1）【同试卷 IV 第二、（1）题】 （2）B。 （3）B。 （4）B。 （5）【同试卷 IV 第二、（4）题】


三、\(\displaystyle \dfrac{\displaystyle 1}{\displaystyle 2}\)。


四、【同试卷 IV 第五题】


五、\(\displaystyle x^2\cos x-4x\sin x-6\cos x+C\)，其中 \(\displaystyle C\) 为任意常数。


六、甲种鱼放养数为 \(\displaystyle \dfrac{\displaystyle 3\alpha-2\beta}{\displaystyle 2\alpha^2-\beta^2}\)，乙种鱼放养数为 \(\displaystyle \dfrac{\displaystyle 4\alpha-3\beta}{\displaystyle 2(2\alpha^2-\beta^2)}\) 时，产鱼总量最大。


七、（1）\(\displaystyle a=\dfrac{\displaystyle 1}{\displaystyle e}\)，切点为 \(\displaystyle (e^2,\allowbreak 1)\)。（2）\(\displaystyle S=\dfrac{\displaystyle 1}{\displaystyle 6}e^2-\dfrac{\displaystyle 1}{\displaystyle 2}\)。


八、【同试卷 IV 第六题】


九、证明略。（先证明 \(\displaystyle \alpha_1+\alpha_2,\allowbreak \alpha_2+\alpha_3,\allowbreak \alpha_3+\alpha_1\) 都是 \(\displaystyle Ax=0\) 的解，再证明 \(\displaystyle k_1(\alpha_1+\alpha_2)+k_2(\alpha_2+\alpha_3)+k_3(\alpha_3+\alpha_1)=0\Rightarrow k_1=k_2=k_3=0\)。）


十、【同试卷 IV 第十题】


十一、\(\displaystyle P\{V_n=m\}=\dbinom{\displaystyle n}{\displaystyle m} 0.01^m0.99^{n-m}\)（\(\displaystyle m=0,\allowbreak 1,\allowbreak 2,\allowbreak \ldots,\allowbreak n\)）。


十二、【同试卷 IV 第十二题】

\end{document}
